\chapter{Introduction}

\section{Background}
As food insecurity increases and available farmland decreases due to deforestation and urbanization \citep{Olagunju2015}, small-scale greenhouses offer a practical way to grow food in limited spaces. Greenhouses help protect crops from harsh weather and allow for better control of growing conditions, which can improve yield, product quality and lengthen market availability \citep{vox}. By using transparent materials to trap heat and light, they create a more stable environment for plants. However, maintaining this environment can be labour intensive and often requires regular monitoring to manage temperature, humidity, and watering.

\section{Objectives}

\begin{itemize}
    \item Design a small-scale greenhouse that can monitor and adjust key environmental factors such as temperature, soil moisture, and humidity.
    \item Allow users to input and control target conditions based on the specific requirements of different plant types.
    \item Integrate a solar power system to enable off-grid operation.
    \item Address limitations in simultaneous parameter control by prioritizing adjustments that best support overall plant health.
    \item Test and evaluate the performance of the greenhouse system to assess its effectiveness and reliability in maintaining suitable growing conditions.
\end{itemize}

\section{Motivation}

Greenhouses provide ideal growing conditions but require constant monitoring and adjustments, making them labour-intensive and inaccessible for many. By reducing manual effort and enabling customization for different plants, this project aims to empower users to grow food efficiently and sustainably with minimal effort. Additionally, the off-grid solar design promotes energy independence and environmental sustainability, aligning with global efforts to combat food insecurity and climate change. This project bridges technology and agriculture, making greenhouse farming accessible, efficient, and eco-friendly.

\section{Methodology}

This project adopts the engineering design methodology outlined by Dieter and Schmidt, providing a structured framework to guide the development of the electronic greenhouse. \citep{Dieter} This methodology encompasses several key phases:

\begin{enumerate}
    \item \textbf{Problem Definition and Need Identification:} This initial phase emphasizes understanding the problem's scope and identifying the specific needs the design aims to fulfill.
    
    \item \textbf{Concept Generation:} Developing multiple design concepts that address the identified needs and constraints. Multiple potential solutions are brainstormed without immediate judgment, encouraging a broad exploration of ideas.
    
    \item \textbf{Concept Selection:} Here, the various concepts are evaluated against predefined criteria to select the most promising solution.
    
    \item \textbf{Detailed Design:} Precise specifications for all components are established, including materials, dimensions, and tolerances, preparing the design for manufacturing.
    
    \item \textbf{Design Evaluation:} Testing the designed system to determine its functionality and performance.
\end{enumerate} By following this systematic approach, the project ensures a thorough exploration of design alternatives and a robust development process for the electronic greenhouse.